\subsection{Linear topologies}

The category of abelian groups can be upgraded to a symmetric monoidal category $$(\textbf{Ab}, \otimes, \mathbb{Z}).$$
A natural question is if this construction can be lifted to the category $\textbf{TopAb}$. A hopeful first step to answering this question is the fact that $\textbf{TopAb}$ has tensor products, namely, the tensor product $G \otimes H$ in $\textbf{Ab}$ equipped with the finest topology which makes the natural map
\[G \times H \rightarrow G \otimes H\]
continuous satisfies the expected universal property (i.e. that any continuous bilinear map $H \times H \rightarrow Z$ factors uniquely through the tensor product). The natural candidate for the identity would be $\mathbb{Z}$ with the discrete topology. The question now becomes whether or not $(\textbf{TopAb}, \otimes, \mathbb{Z})$ is \emph{not} a monoidal category. We would like to start this section by explaining why this does \emph{not} give a monoidal category, which will thus motivate the introduction of linearly topologized abelian groups later on.

\[\text{\TODO explain it}\]

\begin{definition} A TA $G$ is called a linearly topologized abelian group (LTA for short) if it admits an FSN $\mathcal{U}$ given by open subgroups, which we call an FSG (fundamental system of subgroups).
\end{definition}

\begin{remark} Note that if $\mathcal{U}$ is a pre-FSN of $G$ composed of subgroups, then $\mathcal{U}$ will be a FSG for the topology generated by $\mathcal{U}$. We call such a $\mathcal{U}$ a pre-FSG.
\end{remark}

\noindent Let $\textbf{LTopAb}$ denote the full subcategory of $\textbf{TopAb}$ whose objects are LTA's. The inclusion admits a left-adjoint
\begin{equation}
    L : \textbf{TopAb} \leftrightharpoons \textbf{LTopAb} : \; \supseteq,
\end{equation}
where the functor $L$ sends a TA $G$ to the same group, but with the coarser topology generated only by the open subgroups of $G$. This means $\textbf{LTopAb}$ is a \emph{reflective} subcategory of $\textbf{TopAb}$, and therefore all limits and colimits that exist in the second also exist in the first. More concretely, any limit of LTA's in the bigger category is guaranteed to be an LTA, and any colimit can be transfered via the left-adjoint $L$.

\begin{definition} We say a morphism $f : G_1 \times G_2 \rightarrow H$ of LTA's is \emph{strongly continuous} if for any open subgroup $V \subseteq H$ there are open subgroups $U_1 \subseteq G_1$ and $U_2 \subseteq G_2$ such that \[f(G_1, U_2) + f(U_1, G_2) \subseteq V.\]
\end{definition}

\begin{definition} Given LTA's $G$ and $H$, define their \emph{strong tensor product} $G \circledast H$ as the tensor product $G \otimes H$ in $\textbf{Ab}$ equipped with the topology generated by the pre-FSG given by the image of $G \otimes V + U \otimes H$ in $G \otimes H$ for all open subgroups $V \subseteq H$ and $U \subseteq G$.
\end{definition}

\begin{proposition} Let $G$ and $H$ be LTA's. The natural bilinear map $G \times H \rightarrow G \circledast H$ is strongly continuous and universal in the following sense: for any other bilinear strongly continuous morphism $G \times H \rightarrow Z$, there is a unique continous morphism $G \circledast H \rightarrow Z$ making the obvious diagram commute.
\end{proposition}
\begin{proof} \TODO
\end{proof}


\noindent Note that $G \circledast H$ and $G \otimes H$ are equal as groups, but the topology on the first one is \emph{coarser}. It is precisely this change which will allow us to upgrade $\textbf{LTopAb}$ into a symmetric monoidal category, as stated below.

\begin{theorem}\label{thm:ltopab_is_monoidal} The category $(\textbf{LTopAb}, \circledast, \mathbb{Z})$ is symmetric monoidal.
\end{theorem}
\begin{proof} \TODO
\end{proof}
